\section{Fixed-\texorpdfstring{$p$}{p} bit complexity}\label{sec5}\label{sec:complexity}

This section turns the geometry and the decision procedure into a \emph{bit complexity} statement: for
fixed output dimension $p$ and a \emph{geometrically simple real spectrum}, deciding feasibility in the
hyperplane layer takes time polynomial in the bit size of the input. This is the only complexity result
of the paper and the only one carrying a promise; its boundaries are in
Section~\ref{sec5:sub:scope}. Throughout, $\tau$ is the encoding length of $(A,C,L_0,q)$, all
arithmetic is exact, and proofs of the estimates are in Appendix~\ref{appBC}, where the explicit
bounds are recorded.

\subsection{Geometric simplicity}\label{sec5:sub:gsimple}

\begin{defn}[Geometrically simple real spectrum]\label{sec5:def:gsimple}
Let $M=\pi_{K_0}A|_{K_0}$ be the compression. We say that $M$ has a \emph{geometrically simple real
spectrum} if $\dim\ker(M-\mu I)=1$ for every real eigenvalue $\mu$ of $M$.
\end{defn}

Each real eigenvalue then contributes exactly \emph{one} projective direction. The condition is weaker
than algebraic simplicity: an eigenvalue may have algebraic multiplicity $>1$ while its eigenspace
stays one-dimensional, and it still contributes a single direction, the base
$\mathbb{P}(E_\mu)$ being a point. Conversely, if some real eigenvalue has
$d_\mu=\dim E_\mu\ge2$, the base has dimension $d_\mu-1$ and the number of parameters grows with $n$
(Section~\ref{sec5:sub:scope}). The prescribed rows enter through
$E_\mu^0=E_\mu\cap\ker(L_0N)\cap\ker(L_0AN)$: the enumerable directions are those of $E_\mu^0$, so
under geometric simplicity each real eigenvalue yields at most one candidate, and an eigenvalue with
$E_\mu^0=\{0\}$ is excluded outright.

\subsection{Arithmetic model for the number field}\label{sec5:sub:model}

All exact computations of the decision procedure take place in the number field generated by an
eigenvalue, and the bit bounds below refer to a fixed representation of that field. The representation
is chosen once and for all.

\begin{defn}[Effective representation of $K$]\label{sec5:def:model}
Let $\mu$ be a real eigenvalue of $M$ and let $f_K\in\mathbb{Z}[\theta]$ be the minimal polynomial of
$\mu$ over $\mathbb{Q}$, taken primitive with positive leading coefficient. Put
$d_K=\deg f_K\le k\le n$ and let $\beta_K$ be the largest bit length of the coefficients of $f_K$. The
field is $K=\mathbb{Q}(\mu)\cong\mathbb{Q}[\theta]/(f_K)$, and an element of $K$ is represented by its
coordinate vector in the power basis $(1,\theta,\dots,\theta^{d_K-1})$, each coordinate a rational.
The \emph{size} of an element is the largest bit length of its coordinates and is written $\beta$.
\end{defn}

Three points make this representation usable in a bit model.

\emph{(i) Field parameters.} By Lemma~\ref{sec5:lem:height}(c) below, $d_K\le n$ and
$\beta_K=\poly(n,\tau)$: $f_K$ is a factor of the integer polynomial $\chi_M$, and the height of a
factor is bounded by the height and the degree of the polynomial it divides.

\emph{(ii) Cost of arithmetic.} Addition and multiplication are polynomial arithmetic in
$\mathbb{Q}[\theta]$ followed by reduction modulo $f_K$: $O(d_K^{2})$ coefficient products, each on
rationals of bit length $\poly(d_K,\beta_K)$, hence bit cost $\poly(d_K,\beta_K)$. Inversion uses the
extended Euclidean algorithm in $\mathbb{Q}[\theta]$, of the same order. The zero test is the exact
test on coordinates; no tolerance or numerical approximation enters.

\emph{(iii) The real embedding.} The element $\mu$ is a real root of $\chi_M$, isolated by exact real
root isolation (S2) into a rational interval $I$ that contains no other root. The embedding
$\sigma\colon K\to\mathbb{R}$ sends $\theta$ to that root; it is used only to decide \emph{which} real
root of $f_K$ a direction corresponds to, never to evaluate arithmetic, which stays in $K$. If $f_K$
has several real roots, they are Galois-conjugate and the corresponding instances are conjugate as well,
so the verdict is the same for all of them and the choice of $I$ does not affect it.


\subsection{The bounds}\label{sec5:sub:height}\label{sec5:sub:charpoly}\label{sec5:sub:eigvec}

\begin{lem}[Heights of the data, of the roots and of the field]\label{sec5:lem:height}\label{sec5:lem:charpoly}\label{sec5:lem:roots}\label{sec5:lem:eigvec}
Let the entries of the input have bit length at most $\tau$.
\begin{enumerate}[label=(\alph*),leftmargin=*]
\item \emph{Compression.} A primitive integer basis $N$ of $K_0$ and the compression
$M=(N^{\mathsf T}N)^{-1}N^{\mathsf T}AN$ can be computed in bit complexity $\poly(n,\tau)$, and the
entries of both have bit length $\poly(n,\tau)$.
\item \emph{Characteristic polynomial.} After clearing denominators of $\chi_M(t)=\det(tI-M)$ and
taking the primitive part, the coefficients are integers of bit length $\poly(n,\tau)$, computable in
bit complexity $\poly(n,\tau)$.
\item \emph{Roots and field parameters.} Every root $\mu$ of $\chi_M$ satisfies
$[\mathbb{Q}(\mu):\mathbb{Q}]\le k\le n$ and $h(\mu)\le\poly(n,\tau)$, where $h$ is the logarithmic
absolute height; the real roots can be isolated into rational intervals of any prescribed accuracy in
bit complexity $\poly(n,\tau)$. The defining polynomial of Definition~\ref{sec5:def:model} satisfies
$d_K\le k\le n$ and $\beta_K=\poly(n,\tau)$, and is obtained in bit complexity $\poly(n,\tau)$.
\item \emph{Eigenvectors and fibre bases.} Under Definition~\ref{sec5:def:gsimple}, for every real
eigenvalue $\mu$ the space $E_\mu$ has a primitive vector whose coordinates in $K$ have size
$\poly(n,\tau)$, and bases of $E_\mu^0$, $U_v^\perp$, $D_v=K_0\cap U_v^\perp$ and
$\mathcal{Q}_v/\overline{V}_{0,v}$ can be computed over the same field $K$, in bit complexity
$\poly(n,\tau)$, with element sizes $\poly(n,\tau)$.
\end{enumerate}
\end{lem}

\begin{proof}
(a) Clear denominators, run fraction-free (Bareiss) elimination, normalize primitively and compress;
Hadamard's inequality controls every intermediate minor. (b) The coefficients are signed sums of
principal minors. (c) The degree is at most $\deg\chi_M=k$, and Cauchy's bound turns the coefficient
size of $\chi_M$ into a bound on $h(\mu)$; real root isolation is polynomial-time for integer
polynomials. The factor $f_K$ of the integer polynomial $\chi_M$ has height bounded by the height and
degree of $\chi_M$, which gives $\beta_K=\poly(n,\tau)$. (d) Each of the four spaces is a kernel, image,
sum or quotient of matrices of size at most $n$ over $K$, so fraction-free elimination over $K$
produces its basis, the coordinates being ratios of minors. The explicit bounds and their derivations
are in Appendix~\ref{appBC}, Lemmas~\ref{appBC:lem:norm}--\ref{appBC:lem:bases}.
\end{proof}

\begin{remark}[the bases must be computed over $K$]\label{sec5:rem:sameK}
These bases must be computed over $K=\mathbb{Q}(\mu)$, not over $\mathbb{Q}$: when $\mu\notin\mathbb{Q}$
the space $E_\mu$ is a line over $K$ but its trace over $\mathbb{Q}$ is a sum of conjugate lines, so a
computation over $\mathbb{Q}$ would split one direction into $[K:\mathbb{Q}]$ of them. Direction
enumeration and fibre construction therefore take place in the same field.
\end{remark}

\begin{lem}[Size of the charts]\label{sec5:sub:chart}\label{sec5:lem:chart}
Fix a real eigenvalue $\mu$, a direction $[v]\in\mathbb{P}(E_\mu^0)$ and an affine chart of the fibre
$\Gr(s-\delta,\mathcal{Q}_v/\overline{V}_{0,v})$, with $s=q-k+1$,
$\delta=\dim V_0-\dim(V_0\cap K_0)$, $V_0=\row L_0$, and $\widetilde r:=r(v)-\delta$ the rank of the
quotient bundle (unrelated to the free block $R$, which no longer appears). A chart is given by a
symbolic $d_{\mathrm{chart}}\times(\widetilde r-d_{\mathrm{chart}})$ matrix $X$. Then:
\begin{enumerate}[label=(\roman*),leftmargin=*]
\item $d_{\mathrm{chart}}=(s-\delta)(r(v)-s)\le p^{2}$;
\item every chart row is affine in $X$, and every maximal minor has degree $\le\omega$ in $X$ and in
$\lambda$;
\item the number of maximal minors is $\binom{q}{\omega}\binom{k}{\omega}=k\binom{q}{s}=n^{O(p)}$ for
fixed $p$;
\item each fibre is covered by $\binom{\widetilde r}{s-\delta}=O(1)^{p}$ charts;
\item there are at most $k\le n$ eigenvalues, each contributing at most one direction.
\end{enumerate}
\end{lem}

\begin{proof}
Each item is a count: $x\mapsto x(\widetilde r-x)$ is maximal at $x=\widetilde r/2$ with
$\widetilde r\le\rho\le2p$; the chart rows are affine in $X$ and the pencil is linear in $\lambda$;
$\binom{k}{\omega}=k$ because $\omega=k-1$, and $\binom{q}{s}\le\sum_{j\le2p}\binom{n}{j}$ because
$s\le\rho\le2p$. The counts and the dense-size bound $D=n^{O(p^{2})}$ are in Appendix~\ref{appBC},
Lemma~\ref{appBC:lem:chart}.
\end{proof}

\begin{lem}[Cost of the gcd]\label{sec5:sub:subres}\label{sec5:lem:subres}\label{sec5:lem:multigcd}
Let $K$ be a number field in the representation of Definition~\ref{sec5:def:model}, with parameters
$d_K\le n$ and $\beta_K=\poly(n,\tau)$. Let $f,g\in K[X][\lambda]$ with $\deg_\lambda\le\omega$ and
$\deg_X\le\omega$ in each of the $d$ parameters, $d\le p^{2}$, and coefficient size at most $\beta$ (in
the sense of Definition~\ref{sec5:def:model}). Put
\begin{equation}\label{sec5:eq:sizes}
\begin{gathered}
\beta_S:=2\omega\bigl(\beta+\omega\log_2(2\omega)\bigr)+\poly(d_K,\beta_K),\\
D:=\binom{d+2\omega^{2}}{d}.
\end{gathered}
\end{equation}
Then:
\begin{enumerate}[label=(\alph*),leftmargin=*]
\item \emph{Sizes.} $\operatorname{Syl}(f,g)$ is a $(2\omega)\times(2\omega)$ matrix; the chain has at
most $\omega+1$ subresultants $S_0,\dots,S_\omega$; each $S_j$ lies in $K[X][\lambda]$, has
$\deg_\lambda S_j\le j$, $\deg_X S_j\le 2\omega^{2}$, and coefficient size at most $\beta_S$.
\item \emph{Gcd.} The primitive part of the gcd over $K(X)[\lambda]$ is $S_{j^{*}}$ up to a unit of
$K[X]$, where $j^{*}=\max\{j:S_j\neq0\}$; hence it has $\deg_\lambda=j^{*}$ and coefficient size at
most $\beta_S$. In particular the test ``$\deg_\lambda\gcd>0$'' is the test $j^{*}>0$.
\item \emph{Cost.} The chain is computed by the fraction-free subresultant algorithm in bit complexity
$\poly(D,\beta_S)$, and $D=n^{O(p^{2})}$ once $p$ is fixed.
\item \emph{Several polynomials.} For $m_1,\dots,m_\ell\in K[X][\lambda]$ with $\ell=n^{O(p)}$ and the
same degree and size bounds, the iterated gcd $\gcd(\dots\gcd(m_1,m_2),\dots,m_\ell)$ costs
$\poly(\ell,D,\beta_S)$ and has coefficient size at most $\beta_S$ at every step. The size is
\emph{bounded}, not necessarily decreasing: the primitive gcd is a factor of the current input, so it
satisfies the same factor-height bound as that input, and every intermediate subresultant satisfies the
subresultant bound of item (a). The number of non-trivial steps is at most $\omega$, since the
$\lambda$-degree strictly decreases whenever the gcd is non-trivial.
\end{enumerate}
\end{lem}

\begin{proof}
(a) Each entry of the Sylvester matrix is a coefficient of $f$ or $g$, of $X$-degree $\le\omega$ and
size $\beta$; a $(2\omega)\times(2\omega)$ determinant has $X$-degree at most $2\omega^{2}$ and, by
Hadamard's inequality applied to the $2\omega$ rows, coefficient size at most $\beta_S$ as in
\eqref{sec5:eq:sizes}. (b) The subresultant chain satisfies $\deg_\lambda\gcd=\max\{j:S_j\neq0\}$ and
$S_j\in K[X][\lambda]$; the primitive gcd divides the subresultant, hence obeys the same bound. (c) The
dense representation of a polynomial with $\lambda$-degree $\le\omega$ and $X$-degree $\le 2\omega^{2}$
in $d$ variables has $D=\binom{d+2\omega^{2}}{d}$ coefficients, and the algorithm performs $O(\omega)$
pseudo-remainder steps on such objects. (d) By Gauss's lemma over the UFD $K[X]$, the primitive gcd is a
factor of $m_1$ in $K[X][\lambda]$, so by the factor-height bound (S5) its coefficient size is at most
$\beta_S$; the iteration therefore stays within the bound rather than shrinking it. The counts are in
Appendix~\ref{appBC}, Propositions~\ref{appBC:prop:iter} and~\ref{appBC:prop:cost}, and the failure of
the naive Euclidean algorithm is shown there on an explicit example.
\end{proof}

\subsection{Main theorem}\label{sec5:sub:main}

\begin{thm}[Fixed $p$, geometrically simple spectrum]\label{sec5:thm:poly}
Let $p$ be fixed. On inputs with geometrically simple real spectrum
(Definition~\ref{sec5:def:gsimple}), deciding feasibility in the hyperplane layer ($c=1$) can be
carried out in $\poly(n,\tau)$, concretely $n^{O(p^{2})}$, bit operations, using exact arithmetic
throughout.
\end{thm}

\begin{proof}
The steps and their costs are as follows.
\begin{enumerate}[label=(\arabic*),leftmargin=*]
\item \emph{Compression:} $\poly(n,\tau)$ by Lemma~\ref{sec5:lem:height}(a).
\item \emph{Factorization, field and real roots:} factor $\chi_M$ into irreducibles by (S1), build $K$
as in Definition~\ref{sec5:def:model}, and isolate the real roots by (S2), all $\poly(n,\tau)$; by
Lemma~\ref{sec5:lem:height}(c) each real root $\mu$ satisfies $[K:\mathbb{Q}]\le n$,
$h(\mu)=\poly(n,\tau)$ and $\beta_K=\poly(n,\tau)$.
\item \emph{Directions:} under geometric simplicity $E_\mu$ is one-dimensional and $E_\mu^0$ is a
subspace of it, so each $\mu$ yields at most one candidate direction; the eigenvector and the test for
membership in $E_\mu^0$ cost $\poly(n,\tau)$ by Lemma~\ref{sec5:lem:height}(d).
\item \emph{Fibres:} bases of $D_v$, $U_v^\perp$ and $\mathcal{Q}_v/\overline{V}_{0,v}$, again by
Lemma~\ref{sec5:lem:height}(d).
\item \emph{Charts:} by Lemma~\ref{sec5:lem:chart} there are $\le2^{2p}$ charts per fibre,
$n^{O(p)}$ maximal minors, and degrees $\le\omega\le n$, while $d_{\mathrm{chart}}\le p^{2}$ is a
constant because $p$ is fixed.
\item \emph{Decision:} on each chart compute the subresultant gcd by
Lemma~\ref{sec5:lem:subres}(a)--(d), at cost $\poly(D,\beta_S)$ with $D=n^{O(p^{2})}$ and
$\beta_S=\poly(n,\tau)$ by \eqref{sec5:eq:sizes}; $D$ is polynomial in $n$ since $p$ is fixed.
\item \emph{Decision rule:} by the full-rank leading-coefficient lemma of Section~\ref{sec4},
$\rank H\equiv\omega$ on every chart, so ``the gcd over $K(X)[\lambda]$ has degree zero'' is equivalent
to ``every point of the chart is good''. Hence the layer is feasible if some chart is generically
good, and infeasible if every chart is entirely bad.
\end{enumerate}
There are at most $k\cdot2^{2p}$ pieces and each term is either $\poly(n,\tau)$ or $n^{O(p^{2})}$, so
the total is $n^{O(p^{2})}$; the arithmetic in $K$ costs $\poly(d_K,\beta_K)$ per operation, which is
absorbed in the same bound.
\end{proof}

\paragraph{Promise class} The theorem is a \emph{promise-class} result, not a complexity
classification of the general problem: its first step checks Definition~\ref{sec5:def:gsimple}, and if
some real eigenvalue has eigenspace dimension $\ge2$ the algorithm reports \emph{out of scope}, not
\emph{infeasible}. The theorem asserts nothing about inputs outside the promise.

\subsection{The repeated case: what the counterexample does and does not show}\label{sec5:sub:scope}

\paragraph{Why the promise cannot be dropped \emph{for free}} If some real eigenvalue has
$d_\mu=\dim E_\mu\ge2$, the base at $\mu$ has dimension $d_\mu-1$ and the number of chart parameters
becomes $(d_\mu-1)+(s-\delta)(r(v)-s)$; the second term is still $\le p^{2}$, but the first grows with
$n$, so the dense size becomes $n^{O(d_\mu+p^{2})}$, which is not polynomial for fixed $p$. Three
things must be kept apart: the geometry survives (the parameterization holds for arbitrary $d_\mu$),
the exact decision survives (the chartwise computation is correct for arbitrary $d_\mu$), and only the
\emph{complexity bound} fails.

\paragraph{What the example of Section~\ref{sec:counterexample} proves} It proves that inside a
repeated eigenspace the verdict can depend on the projective direction, so that selecting one
eigenvector per eigenvalue is not a valid reduction. It does \emph{not} prove that the repeated case is
computationally hard, and it does not prove that geometric simplicity is necessary for polynomial-time
solvability: no lower bound of any kind is asserted here. The complexity of the general
(repeated-spectrum) case is left open.

\paragraph{Other cases not covered} Layers with $c\ge2$ admit no corresponding direction--fibre
decomposition and are outside the scope of this paper. The design matrix is assumed to have full row
rank, which forces $\rank L_0=r_0$ as a necessary condition for feasibility and is used as a zero-cost
precondition. Real points cause no difficulty, every parameter chart being a rational real chart with
Zariski-dense real points.

\subsection{Standard facts used}\label{sec5:sub:facts}

Five facts are quoted and not reproved: \emph{(S1)} factorization of integer polynomials into
irreducibles in bit complexity $\poly(N,\beta)$ \cite{vzgg03}; \emph{(S2)} real root isolation into
pairwise disjoint rational intervals at the same cost \cite{bpr06}; \emph{(S3)} arithmetic in a number
field $K=\mathbb{Q}[\mu]/(f)$ of degree $d$ and height $h$ in $\poly(d,h)$ bits with an exact zero
test \cite{lan02,vzgg03}; \emph{(S4)} the subresultant chain yields the primitive part of the gcd over
the fraction field, every intermediate quantity being a Sylvester minor with polynomially controlled
bit length \cite{vzgg03}; \emph{(S5)} the height of a factor of a polynomial is bounded by a polynomial
in the degree and the height of that polynomial (the factor-height, or Mignotte, bound), so the gcd of
a family of polynomials is no larger, in the sense of coefficient size, than the polynomials it divides
\cite{vzgg03,lan02}. Only the existence of polynomial-time algorithms is used, not the constants.
